% !TeX encoding = UTF-8
\documentclass[variant=apuntes,cover=institutional,mode=screen]{ua-engineering-v6-article}
% !TeX encoding = UTF-8
\UASetUniversity{Universidad Austral}
\UASetFaculty{Facultad de Ingeniería}
\UASetDepartment{Departamento de Computación}
\UASetCampus{Pilar}
\UASetCareer{Ingeniería Informática}
\UASetCommission{Comisión A}
\UASetProfessor{Equipo docente}
\UASetSemester{Segundo cuatrimestre 2026}
\UASetCourse{Sistemas Distribuidos}
\UASetDate{2026}

\UASetSemester{Segundo cuatrimestre 2026}
\UASetDate{2026}

\UASetClassNumber{07}
\UASetClassTitle{Logs distribuidos y streaming}
\title{Clase 07 --- Kafka, offsets, semánticas de entrega y estado en movimiento}
\author{\UAProfessor}
\date{\UADate}

\begin{document}
\maketitle

\section{Propósito de la clase}

La clase anterior estudió cómo un nodo conserva historia mediante WAL, versiones e índices. Esta clase amplía la idea: un log puede dejar de ser un detalle interno del motor y convertirse en una interfaz distribuida entre productores y consumidores independientes.

El cambio conceptual es importante. En un sistema basado principalmente en RPC, un componente intenta obtener una respuesta inmediata de otro. En un sistema basado en log, el productor registra un hecho durable y los consumidores avanzan a su propio ritmo. Esto desacopla temporalmente, permite replay y facilita nuevos consumidores, pero no elimina los problemas distribuidos: los desplaza hacia particiones, offsets, estado de progreso, semánticas de entrega y consistencia entre el log y los efectos externos.

\begin{uakeybox}
Un log distribuido no “resuelve” automáticamente orden, entrega o exactly-once. Hace esas decisiones observables y configurables dentro de fronteras concretas.
\end{uakeybox}

Al terminar la clase, el estudiante debe poder:
\begin{itemize}
\item distinguir log distribuido, cola y base orientada a estado;
\item explicar topics, particiones, keys, offsets y consumer groups;
\item construir timelines de pérdida y duplicación;
\item definir con precisión el alcance de at-most-once, at-least-once y exactly-once;
\item explicar producer idempotente, transacciones y transactional outbox;
\item razonar sobre streaming stateful, event time, ventanas y eventos tardíos;
\item proponer métricas y SLO para un pipeline real.
\end{itemize}



\section{Cómo estudiar este capítulo}

Este capítulo está organizado en tres pasadas. La primera construye una intuición con un caso real; la segunda formaliza el mecanismo y sus supuestos; la tercera somete la solución a una falla o a una medición. Para estudiar de manera eficaz:

\begin{enumerate}
\item explique primero el problema sin mencionar una tecnología;
\item dibuje la ejecución normal y mueva el punto de falla;
\item identifique la garantía exacta y su frontera;
\item anote qué costo se desplaza hacia latencia, almacenamiento, reparación u operación;
\item cierre proponiendo una métrica o un experimento que podría refutar su diseño.
\end{enumerate}

\begin{uainfo}
Las analogías del capítulo son puertas de entrada, no definiciones finales. Una respuesta completa debe poder abandonar la analogía y utilizar el vocabulario técnico correspondiente.
\end{uainfo}

\section{Mapa intuitivo antes del formalismo}

Antes de introducir la terminología de Kafka conviene fijar siete imágenes mentales. No sustituyen a las definiciones formales; sirven para responder primero \emph{qué problema existe} y recién después estudiar el mecanismo.

\begin{center}
\begin{tabular}{p{0.18\textwidth}p{0.34\textwidth}p{0.36\textwidth}}
\toprule
Concepto & Intuición & Pregunta técnica que abre \\
\midrule
Log & película de hechos que se conserva & ¿qué orden existe y cuánto tiempo se retiene? \\
Partición & carril ordenado dentro de una autopista & ¿qué eventos deben compartir carril? \\
Offset & señalador de lectura de un consumidor & ¿hasta dónde puede reiniciar sin perder o repetir? \\
Consumer group & equipo que reparte carriles & ¿cuánto paralelismo real existe? \\
Lag & deuda de trabajo acumulada & ¿el consumidor alcanza la tasa de ingreso? \\
Exactly-once & frontera que vuelve atómicos progreso y efecto & ¿qué sistemas participan realmente? \\
Event time & reloj del hecho, no de su llegada & ¿cuánto esperar por eventos tardíos? \\
\bottomrule
\end{tabular}
\end{center}

\begin{center}
\begin{tikzpicture}[node distance=7mm,>=Latex,every node/.style={font=\small,align=center}]
\node[draw=UAIngOrange,rounded corners,thick,fill=UAIngPaper,text width=2.4cm] (p) {Productor\\agrega hechos};
\node[draw=UAIngNavy,rounded corners,thick,right=of p,text width=2.4cm] (l) {Log particionado\\retiene historia};
\node[draw=UAIngOrange,rounded corners,thick,right=of l,text width=2.4cm] (c) {Consumidor\\lee y actualiza};
\node[draw=UAIngNavy,rounded corners,thick,right=of c,text width=2.4cm] (e) {Efecto de negocio\\inventario, pago};
\draw[->,very thick] (p)--(l);
\draw[->,very thick,UAIngOrange] (l)--node[above,font=\scriptsize]{offset}(c);
\draw[->,very thick] (c)--(e);
\end{tikzpicture}
\end{center}

\begin{uainfo}
La pregunta de la clase no es ``¿cómo usar Kafka?'', sino ``¿cómo representar hechos, orden, progreso y efectos para que una falla tenga una respuesta definida?''.
\end{uainfo}

\section{Caso conductor: CampusShop durante una venta relámpago}

Durante diez minutos, CampusShop acepta 600.000 pedidos. El checkout debe responder rápido y publicar un hecho durable. Tres consumidores avanzan a ritmos distintos:

\begin{itemize}
\item \textbf{Inventario}: descuenta stock y no puede aplicar dos veces el mismo pedido.
\item \textbf{Facturación}: genera una factura y debe preservar orden por pedido.
\item \textbf{Analítica}: puede atrasarse, releer y corregir agregaciones.
\end{itemize}

La arquitectura inicial es:

\begin{center}
\begin{tikzpicture}[node distance=5mm,>=Latex,every node/.style={font=\small,align=center}]
\node[draw=UAIngOrange,rounded corners,thick,fill=UAIngPaper,text width=2.1cm] (api) {Checkout\\{\scriptsize\texttt{Order}\\\texttt{Confirmed}}};
\node[draw=UAIngNavy,rounded corners,thick,right=of api,text width=2.1cm] (k) {Topic \texttt{orders}\\particiones};
\node[draw=UAIngOrange,rounded corners,thick,above right=3mm and 8mm of k,text width=1.9cm] (i) {Inventario};
\node[draw=UAIngNavy,rounded corners,thick,right=8mm of k,text width=1.9cm] (f) {Facturación};
\node[draw=UAIngOrange,rounded corners,thick,below right=3mm and 8mm of k,text width=1.9cm] (a) {Analítica};
\draw[->,very thick] (api)--(k);
\draw[->,very thick,UAIngOrange] (k)--(i);
\draw[->,very thick] (k)--(f);
\draw[->,very thick,UAIngOrange] (k)--(a);
\end{tikzpicture}
\end{center}

Este caso permite separar cuatro dimensiones que suelen confundirse:

\begin{enumerate}
\item \emph{durabilidad del registro}: el hecho permanece en el log;
\item \emph{progreso}: el consumidor recuerda hasta dónde avanzó;
\item \emph{estado derivado}: inventario o factura reflejan el hecho;
\item \emph{efecto externo}: un cobro o email ocurre fuera del broker.
\end{enumerate}

\section{Del WAL al commit log distribuido}

WAL y commit log comparten una intuición: conservar una secuencia ordenada de hechos facilita recovery y reproducción. Sin embargo, tienen propósitos distintos.

\begin{itemize}
\item El WAL local suele ser una estructura interna para reconstruir páginas o estado después de un crash.
\item Un log distribuido es una interfaz de integración: múltiples productores escriben y múltiples consumidores leen una historia compartida.
\item El WAL puede compactarse o reciclarse cuando ya no es necesario para recovery; un log distribuido conserva eventos según una política de retención y uso de negocio.
\end{itemize}

La conexión es pedagógicamente valiosa: el offset del consumidor cumple un papel parecido al punto de recuperación. Ambos describen “hasta dónde se ha incorporado la historia”, pero el offset no demuestra que un efecto de negocio externo se haya vuelto durable.

\section{Caso conductor: pedidos durante una campaña masiva}

Supongamos un marketplace que recibe 600.000 eventos de pedido en diez minutos. Los eventos deben alimentar:
\begin{itemize}
\item inventario;
\item facturación;
\item notificaciones;
\item detección de fraude;
\item analytics.
\end{itemize}

Si el checkout invoca sincrónicamente a todos los consumidores, la latencia y disponibilidad del camino crítico quedan acopladas al componente más lento. Un log permite registrar el pedido una vez y que cada consumidor avance con independencia.

Pero aparecen preguntas nuevas:
\begin{itemize}
\item ¿qué orden debe conservarse?;
\item ¿cómo se reparte el trabajo?;
\item ¿desde dónde reinicia un consumidor?;
\item ¿qué ocurre si procesa y falla antes de guardar progreso?;
\item ¿cómo se evita que base de datos y evento diverjan?;
\item ¿qué significa exactly-once para un email o un cobro externo?
\end{itemize}

\section{Commit log, topic y partición}

Un topic es una categoría lógica. Cada partición del topic es una secuencia append-only:
\[
P_j=(e_{j,0},e_{j,1},\ldots)
\]
Cada registro tiene un offset creciente dentro de esa partición.

\subsection{Orden local}

Kafka garantiza orden dentro de una partición. No existe un orden total gratuito entre particiones. Si dos eventos relacionados deben observarse en orden, el producer debe elegir una key que los dirija consistentemente a la misma partición.

Ejemplo: todos los eventos de un pedido pueden usar \texttt{order\_id}. Los tipos\\ \texttt{OrderCreated}, \texttt{PaymentAuthorized} y \texttt{OrderCancelled} mantienen así un orden relativo. Si se usan keys diferentes, los consumidores deben reconstruir causalidad mediante identificadores y reglas adicionales.

\subsection{La key también distribuye carga}

La misma key que preserva orden determina partición y, por lo tanto, distribución. Una key demasiado agregada, como \texttt{country}, puede concentrar millones de eventos de un país en una sola partición. Una key demasiado fragmentada mejora balance, pero puede romper orden semántico.

El diseño correcto declara:
\begin{itemize}
\item unidad de orden;
\item cardinalidad de la key;
\item skew esperado;
\item hot keys plausibles;
\item estrategia para evolucionar el número de particiones.
\end{itemize}

\section{Consumer groups y paralelismo}

Un consumer group representa una aplicación lógica. Las particiones se distribuyen entre sus consumidores activos. Dentro de un grupo, una partición es procesada por un consumidor a la vez, aunque la asignación puede cambiar.

Consecuencias:
\begin{itemize}
\item el máximo paralelismo útil está acotado por la cantidad de particiones;
\item agregar consumidores por encima de ese número deja miembros ociosos;
\item aumentar particiones eleva metadata, conexiones, archivos y complejidad de asignación;
\item una partición caliente puede limitar todo el grupo aunque el promedio parezca sano.
\end{itemize}

\subsection{Rebalance}

Cuando entra o sale un consumidor, el grupo puede reasignar particiones. Durante un rebalance, parte del trabajo se pausa o cambia de dueño. Esto exige:
\begin{itemize}
\item cerrar correctamente recursos;
\item guardar o reconstruir estado local;
\item coordinar offsets;
\item evitar que una partición sea procesada simultáneamente por dos instancias durante una transición insegura.
\end{itemize}

El escalado reactivo demasiado agresivo puede causar una sucesión de rebalances y empeorar el lag que intentaba reducir.

\section{Offsets: progreso durable, no efecto durable}

Conviene distinguir:
\begin{itemize}
\item \textbf{consumer position}: próximo offset que se entregará en la instancia actual;
\item \textbf{committed offset}: posición almacenada para reinicio del grupo.
\end{itemize}

El committed offset suele representar el próximo registro pendiente. Si se confirma 43, el grupo declara que 0--42 fueron incorporados según su política. Pero esa declaración puede ser incorrecta si el efecto downstream y el offset no se coordinan.

\subsection{At-most-once}

Secuencia:
\begin{enumerate}
\item leer offset 42;
\item confirmar 43;
\item fallar antes del efecto.
\end{enumerate}
Al reiniciar, el consumidor continúa en 43. El evento 42 no se reprocesa y su efecto se pierde.

\subsection{At-least-once}

Secuencia:
\begin{enumerate}
\item leer offset 42;
\item aplicar el efecto;
\item fallar antes de confirmar 43;
\item reiniciar en 42 y aplicar de nuevo.
\end{enumerate}
El registro queda disponible para redelivery y el efecto puede duplicarse. La política no garantiza por sí sola que el efecto de negocio llegue a completarse: eso depende de reintentos, manejo de errores permanentes e idempotencia.

\subsection{Lag}

Una definición simple es:
\[
lag=log\_end\_offset-committed\_offset
\]
El lag muestra distancia respecto del final del log. Debe analizarse por partición y junto con tasa de entrada, throughput de consumo y latencia de efecto. Un promedio bajo puede ocultar una partición crítica muy atrasada.

\section{Replicación y ACK del producer}

Cada partición tiene un líder y réplicas. El producer envía al líder y espera un ACK según la configuración.

Las opciones implican distintos trade-offs:
\begin{itemize}
\item esperar poco reduce latencia, pero puede confirmar antes de que otras réplicas tengan el dato;
\item esperar réplicas in-sync mejora durabilidad frente a failover, pero aumenta costo y puede reducir disponibilidad;
\item exigir demasiadas réplicas disponibles puede detener escrituras durante una degradación.
\end{itemize}

La palabra “publicado” debe acompañarse de una semántica: ¿aceptado por el líder?, ¿replicado a ISR?, ¿visible para consumidores read-committed?, ¿durable según el objetivo?

\section{Producer idempotente}

Un producer puede recibir timeout después de que el broker aceptó un batch. Si reenvía, el log podría contener duplicados. Kafka permite que el broker reconozca secuencias de un producer y descarte reenvíos dentro de la sesión y protocolo esperado.

Esto resuelve un tipo específico de duplicación: retry de producción hacia Kafka. No vuelve idempotente:
\begin{itemize}
\item una API bancaria llamada por un consumidor;
\item un correo enviado a un proveedor externo;
\item una escritura a otra base sin coordinación;
\item una acción física.
\end{itemize}

\section{Semánticas de entrega y procesamiento}

Las etiquetas at-most-once, at-least-once y exactly-once son útiles solo si se declara el objeto de la garantía.

\subsection{Entrega física}
¿Cuántas veces recibe el consumidor un registro?

\subsection{Actualización de estado}
¿Cuántas veces cambia el state store o tabla local?

\subsection{Efecto de negocio}
¿Cuántas veces se cobra, descuenta stock, envía una notificación o registra una auditoría?

Un pipeline puede tener exactly-once en Kafka Streams y at-least-once respecto de un sistema externo. Por eso, una arquitectura seria escribe la frontera de transacción.

\section{Exactly-once en un pipeline Kafka--Kafka}

En un pipeline cerrado sobre Kafka, una transacción puede agrupar:
\begin{itemize}
\item lectura y offset de input;
\item actualización del state store con changelog;
\item producción a topics de salida.
\end{itemize}

Si la transacción aborta, el avance y la salida no se hacen visibles parcialmente a consumidores configurados para leer solo registros confirmados. Esta garantía es valiosa, pero no universal. La incorporación de una base externa, un REST API o un correo exige idempotencia, deduplicación o un protocolo adicional.

\section{Dual write}

Un servicio guarda un pedido en una base y publica \texttt{OrderCreated}. Si hace ambas acciones por separado, existe una ventana de falla:
\begin{itemize}
\item DB confirmada y evento ausente;
\item evento publicado y DB abortada;
\item ambos realizados, respuesta perdida y retry.
\end{itemize}

Este problema no se arregla cambiando el orden de las operaciones. Se necesita una forma de hacer atómica la intención o tolerar/reparar la divergencia.

\section{Transactional outbox}

El patrón outbox guarda, en una misma transacción local:
\begin{itemize}
\item la modificación de negocio;
\item un registro que representa el evento pendiente de publicar.
\end{itemize}

Un relay o proceso de CDC publica la outbox al broker. Puede reenviar, por lo que los consumidores deben tolerar duplicados. El patrón elimina la ventana “base sí, evento no” sin exigir una transacción distribuida entre base y broker.

El costo incluye:
\begin{itemize}
\item tabla y retención de outbox;
\item relay/CDC y monitoreo;
\item orden por agregado;
\item deduplicación downstream;
\item latencia adicional antes de publicar.
\end{itemize}

\section{Streaming stateless y stateful}

\subsection{Stateless}
Cada evento se procesa independientemente: filtrar, transformar formato, validar o enriquecer con una tabla inmutable.

\subsection{Stateful}
El resultado depende de eventos previos: agregaciones, joins, sesiones, ventanas, detección de fraude y deduplicación.

El stateful streaming necesita coordinar:
\begin{itemize}
\item estado local;
\item changelog o snapshot;
\item offset de input;
\item output producido;
\item restauración después de fallo o rebalance.
\end{itemize}

\section{Event time, processing time y eventos tardíos}

El tiempo del evento es cuando ocurrió el hecho en el dominio. El tiempo de procesamiento es cuando el sistema lo observó. Pueden separarse por redes móviles, buffering, reintentos o relojes imperfectos.

Para ventanas de negocio, event time suele ser más significativo. Sin embargo, el sistema no sabe con certeza si ya llegaron todos los eventos de una ventana.

\subsection{Watermarks}

Un watermark estima cuánto event time se considera completo. Permite cerrar ventanas con una tolerancia de retraso. Esperar más mejora completitud y aumenta latencia/estado; cerrar temprano produce resultados rápidos y requiere correcciones tardías o descarte.

\subsection{Ventanas}

\begin{itemize}
\item tumbling: intervalos contiguos sin solapamiento;
\item sliding: intervalos solapados que avanzan por paso;
\item session: actividad agrupada hasta un período de inactividad.
\end{itemize}

La elección debe provenir de la semántica del dominio, no del framework.

\section{Caso realista: detección de fraude}

Regla: alertar si una tarjeta aparece en cinco países en diez minutos.

Decisiones:
\begin{itemize}
\item key por tarjeta para mantener estado y orden relativo;
\item ventana por event time;
\item tolerancia a eventos tardíos;
\item estado con TTL;
\item política ante duplicados;
\item manejo de hot keys para tarjetas de alto volumen o ataques.
\end{itemize}

Una alerta temprana puede ser corregida; esperar todos los eventos puede ser demasiado lento. El trade-off no es solo técnico: involucra costo de falso positivo y falso negativo.

\section{Operación del pipeline}

Métricas relevantes:
\begin{itemize}
\item tasa y error del producer;
\item latencia de ACK y retries;
\item bytes/eventos por partición;
\item réplicas fuera de sincronía;
\item lag por partición y edad del evento más antiguo;
\item frecuencia y duración de rebalances;
\item latencia de procesamiento y de extremo a extremo;
\item tamaño del estado y tiempo de restore;
\item eventos tardíos, descartados y corregidos;
\item tasa de duplicados detectados.
\end{itemize}

Un SLO de streaming debe orientarse al usuario o al downstream, por ejemplo: “99.9\% de eventos de pedido visibles en inventario en menos de 5 segundos”, no solo “broker disponible”.

\section{Método de diseño}

Para un pipeline:
\begin{enumerate}
\item definir hechos y contratos de eventos;
\item elegir unidad de orden y key;
\item estimar tasa, skew y crecimiento;
\item elegir particiones y consumer groups;
\item declarar semántica de entrega por consumidor;
\item resolver dual writes;
\item definir estado, tiempo y eventos tardíos;
\item instrumentar lag, error y latencia extremo a extremo;
\item ejecutar fallas y rebalances controlados;
\item documentar límite de exactly-once.
\end{enumerate}


\section{Precisiones técnicas indispensables}

\subsection{Topic, partición, offset e identidad no son sinónimos}

Un topic agrupa una categoría de registros; una partición es una secuencia ordenada; un offset es una posición dentro de esa secuencia. La identidad de negocio debe viajar en el evento porque el mismo hecho puede reaparecer en otro offset después de un retry, una republicación o una migración.

\subsection{Position y committed offset describen progresos diferentes}

La \emph{position} es el próximo registro que una instancia viva intentará leer. El \emph{committed offset} es el punto durable desde el cual el grupo reiniciará. Entre ambos puede existir trabajo entregado pero todavía no confirmado.

\subsection{At-least-once garantiza reintento, no éxito del negocio}

At-least-once evita confirmar progreso antes de tiempo y permite redelivery. No garantiza que un pago, un correo o una escritura externa finalmente tengan éxito: pueden existir errores permanentes, poison messages o dependencias indisponibles. La garantía exacta debe describir qué se reentrega y qué condición hace avanzar el offset.

\subsection{Producer idempotente y consumidor idempotente cubren fronteras distintas}

La idempotencia del producer evita duplicados introducidos por retries de publicación dentro de la sesión y protocolo soportados por Kafka. No evita que un consumidor cobre dos veces. Para eso, efecto de negocio y marca de deduplicación deben persistirse como una unidad atómica en el sistema que posee el efecto.

\subsection{Exactly-once siempre tiene una frontera}

En un pipeline Kafka--Kafka, una transacción puede coordinar offsets, state store con changelog y topics de salida. Un proveedor de email, una API bancaria o una base externa no participan automáticamente. Fuera de la frontera se necesitan idempotencia, deduplicación, outbox/inbox o un protocolo adicional.

\subsection{Outbox elimina una ventana, pero conserva redelivery}

Estado de negocio y evento pendiente se guardan en una transacción local. El relay puede publicar y perder el ACK, por lo que repetirá. El consumidor debe registrar el \texttt{event\_id} en la misma transacción que su efecto local; de otro modo existe una nueva ventana entre efecto y deduplicación.

\subsection{Estado y tiempo también deben recuperarse}

Un operador stateful necesita reconstruir estado, progreso y outputs de forma compatible. En event time, el watermark no descubre el futuro: expresa cuánto retraso se decide esperar. La política de late data debe declarar si descarta, corrige o reabre resultados y qué costo tiene cada opción.

\section{Errores frecuentes}

\begin{itemize}
\item creer que un topic tiene orden global;
\item tratar offset como ID de negocio;
\item confirmar offset sin coordinar efecto;
\item decir “exactly-once” sin frontera;
\item ignorar rebalances y estado local;
\item aumentar particiones sin analizar key y overhead;
\item usar eventos para una interacción que requiere respuesta inmediata sin diseñar el feedback;
\item publicar después de una transacción local sin outbox o reparación;
\item medir solo disponibilidad del broker e ignorar lag y edad.
\end{itemize}

\section{Bibliografía guiada y casos reales}

\begin{itemize}
\item Kleppmann, \emph{Designing Data-Intensive Applications}, capítulos 11 y 12: stream processing, logs, eventos y sistemas derivados.
\item Apache Kafka, \href{https://kafka.apache.org/documentation/}{documentación oficial}: topics, particiones y processing guarantees.
\item Apache Kafka, \href{https://kafka.apache.org/41/operations/basic-kafka-operations/}{Basic Kafka Operations}: particiones, offsets y operación.
\item AWS Prescriptive Guidance, \href{https://docs.aws.amazon.com/prescriptive-guidance/latest/cloud-design-patterns/transactional-outbox.html}{Transactional Outbox}: problema de dual write.
\item Kreps, Narkhede y Rao, \emph{Kafka: a Distributed Messaging System for Log Processing}: motivación y diseño original.
\end{itemize}

\section{Conexión con el Trabajo Final Integrador}

Cada grupo debe identificar:
\begin{itemize}
\item qué interacciones conviene desacoplar;
\item topic, evento, key y owner;
\item orden requerido;
\item semántica de entrega;
\item estrategia de outbox o CDC;
\item SLO de lag/extremo a extremo;
\item efecto externo que queda fuera de exactly-once.
\end{itemize}

Una respuesta insuficiente dice “usaremos Kafka”. Una respuesta defendible explica por qué, para qué eventos, con qué partición, qué consumers y qué comportamiento ante crash.


\section{Ejemplos trabajados paso a paso}

\subsection{Dimensionar particiones con margen y skew}

El topic recibe 80.000 eventos/s. Una instancia de consumidor procesa 5.000 eventos/s de forma sostenible. El mínimo aritmético es:
\[
P_{min}=\left\lceil\frac{80.000}{5.000}\right\rceil=16.
\]
Ese número no es todavía una decisión de producción. Si se reserva 30\% de margen:
\[
P_{cap}=\left\lceil\frac{80.000\times1{,}3}{5.000}\right\rceil=21.
\]
Puede elegirse 24 particiones para permitir crecimiento, siempre que el costo de metadata y operación sea aceptable. Sin embargo, si una key concentra 25\% del tráfico, ninguna media global evita que una sola partición reciba 20.000 eventos/s. El diseño debe combinar capacidad agregada con análisis de skew.

\subsection{Timeline de crash: efecto y offset}

Considere el evento 42 y una tabla de inventario. Hay dos órdenes simples:

\begin{center}
\begin{tabular}{p{0.16\textwidth}p{0.34\textwidth}p{0.34\textwidth}}
\toprule
Paso & Commit antes del efecto & Efecto antes del commit \\
\midrule
1 & leer 42 & leer 42 \\
2 & confirmar 43 & descontar stock \\
3 & crash & crash \\
4 & reiniciar en 43 & reiniciar en 42 \\
Resultado & pérdida de efecto & posible doble descuento \\
\bottomrule
\end{tabular}
\end{center}

La forma robusta para una base transaccional local consiste en aplicar el efecto y registrar el \texttt{event\_id} procesado dentro de la misma transacción. Si el evento reaparece, la restricción única detecta el replay y devuelve el resultado previo.

\subsection{Transactional outbox con un esquema mínimo}

El servicio de pedidos debe guardar el pedido y publicar \texttt{OrderConfirmed}. En vez de dos escrituras independientes, usa una transacción local:

\begin{verbatim}
BEGIN;
INSERT INTO orders(order_id, status, total)
VALUES ('9F2A', 'CONFIRMED', 17500);
INSERT INTO outbox(event_id, aggregate_id, event_type,
                   payload, created_at)
VALUES ('E-771', '9F2A', 'OrderConfirmed',
        '{...}', now());
COMMIT;
\end{verbatim}

Un relay lee la outbox, publica y marca progreso. El relay puede publicar dos veces si el broker acepta y el ACK se pierde. Por eso, outbox resuelve atomicidad entre \emph{estado de negocio e intención de publicar}, pero no elimina la necesidad de deduplicación downstream.

\subsection{Ventanas y evento tardío}

Una compra ocurre a las 10:00:40, el teléfono queda sin red y el evento llega a las 10:07. En una ventana tumbling de cinco minutos:

\begin{itemize}
\item con \textbf{processing time}, el evento cae en 10:05--10:10;
\item con \textbf{event time}, pertenece a 10:00--10:05;
\item con watermark de dos minutos, probablemente llega después del cierre y requiere corrección tardía o descarte;
\item con watermark de diez minutos, se incluye, pero la salida definitiva tarda más y el estado vive más tiempo.
\end{itemize}

No existe un watermark ``correcto'' en abstracto. La decisión depende del costo de esperar y del costo de una respuesta incompleta.

\section{Caso complementario: ranking de Asteria Online}

En el mundo virtual medieval Asteria, cada combate emite eventos \texttt{CombatFinished}. El ranking se calcula por ventanas de cinco minutos.

\begin{itemize}
\item Key por \texttt{player\_id}: preserva orden por jugador, pero un jugador extremadamente activo puede ser una hot key.
\item Semántica: at-least-once con deduplicación por \texttt{combat\_id}.
\item Tiempo: event time del servidor autoritativo, no del cliente manipulable.
\item SLO: 99\% de resultados visibles en ranking en menos de 30 s.
\item Corrección: un resultado tardío puede actualizar el ranking; una recompensa económica basada en ranking exige una etapa de cierre más fuerte.
\end{itemize}

Este ejemplo muestra que el mismo stream puede alimentar una vista eventual para mostrar posiciones y un proceso más controlado para otorgar premios.

\section{Lectura operacional: del síntoma a la causa}

Un dashboard muestra: ingreso 80k eventos/s, consumo total 76k/s, lag promedio moderado, P2 con lag 1,8 millones, 12 rebalances/h y p99 de procesamiento de 800 ms. Una lectura causal razonable es:

\begin{enumerate}
\item el sistema tiene déficit agregado de 4k eventos/s, por lo que el backlog crece;
\item P2 concentra más atraso: investigar key skew y costo por evento;
\item los rebalances frecuentes agregan pausas y restauraciones de estado;
\item aumentar consumers puede no ayudar si P2 sigue siendo una sola partición caliente;
\item se debe medir edad del evento más antiguo y latencia end-to-end, no solo número de offsets.
\end{enumerate}

\section{Glosario operativo}

\begin{description}
\item[Registro.] Unidad append-only almacenada en una partición.
\item[Offset.] Posición monotónica dentro de una partición.
\item[Position.] Próximo offset que leerá una instancia viva.
\item[Committed offset.] Punto durable desde el que reiniciará el grupo.
\item[Lag.] Distancia entre el final y el progreso confirmado.
\item[Rebalance.] Reasignación de particiones entre miembros de un grupo.
\item[Watermark.] Estimación de cuánto event time se considera suficientemente completo.
\item[State store.] Estado local/materializado mantenido por un procesador stateful.
\end{description}

\section{Autoevaluación ampliada}

\begin{enumerate}
\item ¿Qué conserva un log que una cola de trabajo clásica puede descartar?
\item ¿Qué orden garantiza una partición y cuál no garantiza?
\item ¿Por qué el offset no es un identificador de negocio?
\item ¿Cómo limita el número de particiones al paralelismo de un grupo?
\item ¿Qué diferencia existe entre position y committed offset?
\item ¿En qué timeline aparece pérdida at-most-once?
\item ¿En qué timeline aparece duplicación at-least-once?
\item ¿Qué información debe persistirse para volver idempotente un efecto?
\item ¿Qué mide lag y qué no prueba?
\item ¿Por qué una hot partition puede ocultarse en el promedio?
\item ¿Qué costo introduce un rebalance en un pipeline stateful?
\item ¿Qué cubre el producer idempotente?
\item ¿Qué queda fuera de una transacción Kafka--Kafka?
\item ¿Qué ventana de falla elimina transactional outbox?
\item ¿Por qué outbox puede producir duplicados?
\item ¿Qué diferencia hay entre event time y processing time?
\item ¿Qué trade-off expresa un watermark?
\item ¿Qué estado debe recuperarse después de un crash en un join stateful?
\item ¿Qué SLO describiría el efecto útil para el negocio?
\item ¿Qué experimento ejecutaría para validar crash, replay y deduplicación?
\end{enumerate}

\section{Síntesis final}

\begin{uakeybox}
El log convierte historia en una interfaz durable y replayable. La partición define el orden, el offset define progreso, la transacción define visibilidad conjunta y la aplicación define el efecto de negocio. El diseño serio declara cada frontera y la valida con lag, latencia y fallas.
\end{uakeybox}

\end{document}
